\subsection{Sensitivity to second-order wave shape}
\label{sec:crest-sensitivity}

Changing crest shape can move an echo across the power threshold even when the vessel follows the same trajectory. To isolate this effect, we paired each linear regular wave with the second-order deep-water Stokes surface \cite{Chen2018StokesHarmonics}. For fundamental phase $\chi=\mathbf{k}\cdot\mathbf{x}-\omega t+\psi$, the bound harmonic gives
\begin{equation}
 \eta_2=a\cos\chi+\tfrac12 ka^2\cos(2\chi).
\end{equation}
We used $\eta_2$ at the buoy's fixed horizontal location and along the instantaneous ray, retaining the same hull motion and visit times. The surface gradient magnitude is bounded by $ka+(ka)^2$. At the largest linear registration-gain cell, the peak slope rises from $ka=0.192$ to $0.205$, so the perturbed case probes the edge of the nominal $ka\le0.20$ screen.

At the reference point, the 5.55 mm harmonic reduces the unstabilized usable fraction from $92.97\%$ to $91.02\%$. The level-pointed fraction stays at $100\%$. At the largest linear registration-gain cell, the 41.5 mm harmonic reduces the number of body-pointed echoes from 2,338 to 1,960 and changes the number meeting tolerance with the midpoint pose from 878 to 879. The per-spoke registration gain consequently falls from 23.76 to 17.59 percentage points. In the pilot record, the number of echoes acquired drops from 96 to 72 of 192 visits, while per-spoke correction still recovers 48 out-of-tolerance positions. Refining the ray sampling from 601 to 1,201 points changes two level-pointed classifications in the strongest cell and none in the other comparisons. The relative prevalence of the two failure modes thus depends on the modeled crest clearance and buoy elevation.
